%!TEX recipe = xelatex
\documentclass[a1,portrait]{a0poster}

% ----------------------------------------------------------
% Packages
% ----------------------------------------------------------
\usepackage{lipsum}
\usepackage{graphicx}
\usepackage{geometry}
\usepackage{tikz}
\usepackage{tcolorbox}
\tcbuselibrary{skins}
\usepackage{titlesec}
\usepackage{multicol}
\usepackage{xeCJK}
\usepackage{hyperref} % Added for clickable email/links

\setCJKmainfont{MS Mincho} % Or MS Mincho if on Windows

% Page geometry
\geometry{
  left=15mm,
  right=15mm,
  top=0mm, % Top is 0 because our header handles the top margin
  bottom=25mm
}

% ----------------------------------------------------------
% ROBUST HEADER DEFINITION
% ----------------------------------------------------------
\newcommand{\TitleBlock}{
\begin{tikzpicture}[remember picture,overlay]
    \node[
        anchor=north, 
        fill=teal!70!white, 
        minimum width=\paperwidth, 
        minimum height=8cm, % Fixed height for the blue bar
        inner ysep=1cm,     % Padding inside the blue bar
        text=white
    ] at (current page.north) {
        % We use a minipage to hold the content centered
        \begin{minipage}{0.95\paperwidth}
            \centering
            
            % 1. The Title
            {\fontsize{60}{72}\selectfont\bfseries Pre-Processing of Plucked String Instrument Data for Machine Learning}\\[0.5cm]
            
            % 2. Name and Affiliation
            {\Huge Juan Pablo Corredor \quad $\cdot$ \quad 東京藝術大学}\\[0.3cm]
            
            % 3. Contact Information (Email / Web)
            {\Large \texttt{juan.corredor@geidai.ac.jp} \quad $\cdot$ \quad \texttt{github.com/juancorredor}}\\[0.8cm]

            % 4. Logos (Centered)
            % Ensure these image paths are correct or use example-image
            \includegraphics[height=2cm]{example-image} 
            \hspace{2cm} 
            \includegraphics[height=2cm]{example-image}
        \end{minipage}
    };
\end{tikzpicture}
}

% ----------------------------------------------------------
% Section style
% ----------------------------------------------------------
\tcbset{
  sectionbox/.style={
    colback=white,
    colframe=blue!50!white,
    boxrule=0.8pt,
    sharp corners,
    enhanced,
    left=6pt,right=6pt,top=6pt,bottom=6pt, % Slightly more padding
    fonttitle=\bfseries\Large,
    title={#1}
  }
}

\titleformat{\section}{\vspace{0.3cm}}{}{0pt}{}

% ----------------------------------------------------------
% Begin Document
% ----------------------------------------------------------
\begin{document}
\pagestyle{empty}

% 1. Render the Header
\TitleBlock

% 2. CRITICAL: Push the main content down
% Since the header is an "overlay" (to cover the edges), we must
% manually push the text down so it doesn't print ON TOP of the header.
% Adjust this value (9cm) if you make the header taller/shorter.
\vspace*{9cm} 

% 3. Start Content
\begin{multicols}{2}

\begin{tcolorbox}[sectionbox={Background}]
Your background text goes here.
\lipsum[1]
\end{tcolorbox}

\begin{tcolorbox}[sectionbox={Objectives}]
Write your research goals here.
\lipsum[2]
\end{tcolorbox}

\begin{tcolorbox}[sectionbox={Methodology}]
Your methodology description goes here.
\begin{center}
    % Using [width=\linewidth] makes it fit the column perfectly
    \includegraphics[width=\linewidth]{example-image}
\end{center}
\lipsum[3]
\end{tcolorbox}

\begin{tcolorbox}[sectionbox={Results}]
\begin{center}
    \includegraphics[width=\linewidth]{example-image}
\end{center}
\lipsum[4]
\end{tcolorbox}

\begin{tcolorbox}[sectionbox={Conclusion}]
\lipsum[5]
\end{tcolorbox}

% ----------------------------------------------------------
% References
% ----------------------------------------------------------
\vspace{0.5cm}
\noindent\textbf{References}\\
\footnotesize
\begin{itemize}
    \item Author A. *Title of paper*, Journal, Year.
    \item Author B. *Another reference*.
\end{itemize}

\end{multicols}

\end{document}